% 369_Vier_Wege_Leptonmassen_En_ch.tex
% Johann Pascher, 20 September 2026
% Inventory: routes to the lepton masses in the FFGFT corpus

\providecommand{\BM}{\textbf{[B]}}
\providecommand{\KM}{\textbf{[K]}}
\providecommand{\SM}{\textbf{[S]}}
\providecommand{\EM}{\textbf{[E]}}

\chapter*{Four Routes to the Lepton Masses: An Inventory}
\addcontentsline{toc}{chapter}{Four Routes to the Lepton Masses}

\begin{abstract}
The FFGFT corpus now contains four independent approaches to the masses of the
charged leptons: the T0 ladder from torus winding numbers (Docs.~006, 338, 352),
the Koide formula with $\theta = 2/9$ from a single $A_5$ matrix element
(Docs.~292, 367), the Galois identities from GF$(3^n)$ orders (Doc.~338), and the
$\phi$-skeleton of the icosahedral redistribution weights (Doc.~368). This document
places the four routes side by side: input, output, accuracy against PDG,
epistemic status. Finding: the routes deliver different \emph{kinds} of statement
(absolute masses, ratios, constraints) with accuracies from $1{.}6\,\%$ down to
$10^{-3}\,\%$. The Koide route is the only parameter-free one and the most precise;
the T0 ladder the only one yielding absolute masses. A fifth route -- the Koide scale $M$ from the T0 ladder -- is Route~1
summed ($0{.}31\,\%$) and shows that $M$ is not an independent scale.
The open edge is the same in all: the fractal-recursive correction from percent level to $10^{-3}$ level is
measured, not derived.
\end{abstract}

\section*{Notation}
\addcontentsline{toc}{section}{Notation}

\paragraph*{Epistemic markers:}\mbox{}\\[2pt]
\begin{tabular}{@{}ll@{}}
\textbf{[B]} & algebraically proved, exact arithmetic \\
\textbf{[K]} & numerically confirmed (approximation, accuracy stated) \\
\textbf{[S]} & structurally observed, derivation open \\
\textbf{[E]} & established external mathematics or physics \\
\end{tabular}

\medskip
\begin{tabular}{@{}lp{11cm}@{}}
\toprule
\textbf{Symbol} & \textbf{Meaning} \\
\midrule
$\xi$ & FFGFT fundamental parameter, $\xi = 4/30000 = 1/7500$ \\
$v$ & electroweak scale, $246$ GeV \\
$r_i,\,p_i$ & winding quantum numbers of lepton $i$ (Doc.~338, Tab.~1) \\
$K_\text{frak}$ & fractal-recursive correction factor $74/75$ (Doc.~338) \\
$\theta$ & Koide phase, $\theta = 2/9$ \\
$A$ & $\langle v_0\mid R_5\,v_e\rangle$, $|A| = \sqrt2/3$ (Doc.~367) \\
$p_0,p_1,p_2$ & icosahedral redistribution weights (Doc.~293) \\
$\phi$ & golden ratio \\
PDG & Particle Data Group 2024: $m_e = 0{.}51099895$, $m_\mu = 105{.}6583755$, $m_\tau = 1776{.}93\pm0{.}09$ MeV \\
\bottomrule
\end{tabular}

\section{The Four Routes}

\subsection{Route 1: T0 Ladder (Winding Numbers)}

\textbf{Documents:} 006, 338, 352. \quad
\textbf{Input:} $\xi$, $v$, quantum numbers $(r_i, p_i)$. \quad
\textbf{Output:} absolute masses.

\[
  m_i^\text{bare} = r_i\,\xi^{p_i}\,v, \qquad
  \frac{m_j}{m_i} = \frac{r_j}{r_i}\,\xi^{\,p_j - p_i} \quad \BM
\]
with $(r_e,p_e) = (4/3,\,3/2)$, $(r_\mu,p_\mu) = (16/5,\,1)$, $(r_\tau,p_\tau) = (25/9,\,2/3)$
from the torus winding numbers $(n_\theta, n_\phi)$ (Doc.~338, Tab.~1).

\begin{center}
\begin{tabular}{@{}lllll@{}}
\toprule
Quantity & Formula & Value & PDG & Dev. \\
\midrule
$m_\mu/m_e$  & $\tfrac{12}{5}\,\xi^{-1/2}$  & $207{.}846$ & $206{.}768$ & $+0{.}52\,\%$ \\
$m_\tau/m_e$ & $\tfrac{25}{12}\,\xi^{-5/6}$ & $3531{.}6$  & $3477{.}4$  & $+1{.}56\,\%$ \\
$m_e$ [MeV]  & $r_e\,\xi^{3/2}\,v$ & $0{.}5050$ & $0{.}5110$ & $-1{.}18\,\%$ \\
$m_\mu$ [MeV]& $r_\mu\,\xi\,v$     & $104{.}96$ & $105{.}66$ & $-0{.}66\,\%$ \\
$m_\tau$ [MeV]& $r_\tau\,\xi^{2/3}\,v$ & $1783{.}4$ & $1776{.}9$ & $+0{.}37\,\%$ \\
\bottomrule
\end{tabular}
\end{center}

\textbf{Status:} formula \BM, quantum numbers from GF$(3^n)$ \KM. The residual
deviation of $0{.}4$--$1{.}6\,\%$ is the fractal-recursive correction (Doc.~352):
it is \emph{measured} as $+39{.}1\,\xi$ ($m_\mu/m_e$) and $+117\,\xi$
($m_\tau/m_e$; Doc.~352: $117{.}4$ with older PDG value $m_\tau=1776{.}86$); its coefficients are not derived \SM.

\subsection{Route 2: Koide Formula with $\theta = 2/9$}

\textbf{Documents:} 292, 367. \quad
\textbf{Input:} one matrix element $A = \langle v_0\mid R_5\,v_e\rangle$. \quad
\textbf{Output:} mass ratios; scale $M$ external.

\[
  \sqrt{m_k} = M\,a_k, \qquad
  a_k = 1 + 3|A|\,\cos\!\bigl(|A|^2 + \tfrac{2\pi k}{3}\bigr),
  \qquad |A| = \frac{\sqrt2}{3} \quad \BM
\]
The modulus of $A$ supplies the coefficient $\sqrt2$, its modulus squared the
phase $\theta = 2/9$ (Doc.~367). No free parameter.

\begin{center}
\begin{tabular}{@{}lllll@{}}
\toprule
Quantity & Formula & Value & PDG & Dev. \\
\midrule
$m_\mu/m_e$  & $(a_2/a_1)^2$ & $206{.}7703$ & $206{.}7683$ & $0{.}0010\,\%$ \\
$m_\tau/m_e$ & $(a_0/a_1)^2$ & $3477{.}47$  & $3477{.}37$  & $0{.}0031\,\%$ \\
$m_\tau$ from $m_e,m_\mu$, $Q=2/3$ & & $1776{.}97$ & $1776{.}93$ & $0{.}0022\,\%$ \\
$Q_\text{Koide}$ & $(\sum m)/(\sum\sqrt m)^2$ & $2/3$ exact & $0{.}66666$ & $10^{-5}$ \\
\bottomrule
\end{tabular}
\end{center}

\textbf{Status:} $\theta = 2/9$ \BM. $\cos(2/9)$ transcendental \EM, hence no
algebraically exact ratios. The deviations lie \emph{within} the PDG uncertainty
of $m_\tau$ ($\pm0{.}005\,\%$).

\subsection{Route 3: Galois Identities}

\textbf{Document:} 338. \quad
\textbf{Input:} orders of GF$(3^n)$. \quad
\textbf{Output:} two constraints, not a complete determination.

\[
  \left(\frac{m_\mu}{m_e}\right)^2 = 43200 = 2^6\cdot3^3\cdot5^2, \qquad
  m_e\,m_\mu = 54\;\text{MeV}^2 = |\text{GF}(3)^*|\cdot|\text{GF}(27)| \quad \KM
\]

\begin{center}
\begin{tabular}{@{}llll@{}}
\toprule
Identity & Prediction & PDG & Dev. \\
\midrule
$(m_\mu/m_e)^2 = 43200$ & $207{.}846$ & $206{.}768$ & $0{.}52\,\%$ \\
$m_e\,m_\mu = 54$ MeV$^2$ & $54$ & $53{.}991$ & $0{.}016\,\%$ \\
\bottomrule
\end{tabular}
\end{center}

\textbf{Status:} both \KM. The first identity is \emph{the same number} as
Route~1 ($207{.}846 = \sqrt{43200}$): the T0 ladder reproduces the Galois identity
exactly. For $m_\mu/m_e$, Routes~1 and~3 are not two routes but one route in two
languages. The second identity is markedly more precise at $0{.}016\,\%$ and enters
$1/\alpha = 3700/27$ via $\alpha = \xi E_0^2/K_\text{frak}$.

\subsection{Route 4: $\phi$-Skeleton of the $p$-Values}

\textbf{Document:} 368. \quad
\textbf{Input:} $p_0,p_1,p_2$ from $A_5$ (Doc.~293). \quad
\textbf{Output:} ratios algebraically, without cosine.

\[
  \frac{p_1}{p_2} = \phi^8, \qquad \frac{p_0}{p_2} = 2\phi^4 \quad \BM
\]

\begin{center}
\begin{tabular}{@{}lllll@{}}
\toprule
Quantity & Formula & Value & PDG & Dev. \\
\midrule
$m_\tau/m_e$ & $74\,\phi^8$ & $3476{.}42$ & $3477{.}37$ & $0{.}027\,\%$ \\
$m_\tau/m_e$ & $74\,\phi^8\,(1+\tfrac{27}{12}\xi)$ & $3477{.}47$ & $3477{.}37$ & $0{.}0029\,\%$ \\
$m_\mu/m_e$  & $30\,\phi^4$ & $205{.}62$ & $206{.}77$ & $0{.}55\,\%$ \\
\bottomrule
\end{tabular}
\end{center}

\textbf{Status:} $\phi$-powers \BM, Galois factors $74$, $30$, $27/12$ observed
\SM. With the $\xi$ correction, $m_\tau/m_e$ reaches the precision of the Koide route.

\section{Synopsis}

\begin{center}
\begin{tabular}{@{}lllllll@{}}
\toprule
Route & Doc. & Input & Output & $m_\mu/m_e$ & $m_\tau/m_e$ & Status \\
\midrule
1 T0 ladder & 006/338/352 & $\xi, v, r_i, p_i$ & absolute & $0{.}52\,\%$ & $1{.}56\,\%$ & \KM, corr.\ \SM \\
2 Koide & 292/367 & $A$ & ratios & $0{.}001\,\%$ & $0{.}003\,\%$ & $\theta$ \BM \\
3 Galois & 338 & GF$(3^n)$ & constraints & $0{.}52\,\%$ & --- & \KM, cause \SM \\
4 $\phi$-skeleton & 368 & $p_j$ & ratios & $0{.}55\,\%$ & $0{.}027\,\%$ & \BM+\SM \\
4b $\phi$-sk.+$\xi$ & 368 & $p_j,\xi$ & ratios & --- & $0{.}003\,\%$ & \KM \\
\bottomrule
\end{tabular}
\end{center}

\paragraph{What each route delivers.}
\begin{itemize}
  \item \textbf{Route 1} is the only one yielding absolute masses -- and the least
    accurate. Its value is the anchoring to $v$ and $\xi$.
  \item \textbf{Route 2} is the only parameter-free one and the most precise. Its
    limit is the scale: $M \approx 313{.}84$ MeV must come from outside. Both ratios
    lie within the PDG uncertainty of $m_\tau$.
  \item \textbf{Route 3} yields no masses but two constraints. The more precise of
    the two ($m_e m_\mu = 54$) is the anchor of the $\alpha$ derivation.
  \item \textbf{Route 4} is the most algebraic: pure $\phi$-powers from the $A_5$
    geometry, no transcendental function. The price is the Galois factors, which
    are observed and not derived.
\end{itemize}

\paragraph{Two accuracy classes.}
Routes 1, 3 and 4a sit at \emph{percent level} ($0{.}5$--$1{.}6\,\%$). Routes 2
and 4b sit at \emph{$10^{-3}$-percent level}, within experimental uncertainty. The
difference between the classes is precisely the fractal-recursive correction:
$+117{.}4\,\xi$ in Route~1, $+\tfrac{27}{12}\xi$ in Route~4, implicit in
$\cos(2/9)$ in Route~2.

\section{The Common Open Edge}

All four routes end at the same point. The percent class needs a
$\xi$-proportional correction to become the $10^{-3}$ class. This correction is
known in three forms:
\begin{itemize}
  \item as a measured coefficient $117{.}4$ (Route~1, Doc.~352),
  \item as a Galois expression $27/12 = |\text{GF}(27)|/(n_{\phi,\mu}\,n_{\theta,e})$
    (Route~4, Doc.~368, $0{.}5\,\%$ agreement),
  \item as whatever the Koide formula contains beyond $74\phi^8$ (Route~2).
\end{itemize}
A first-principles derivation is missing in all three forms \SM. The proof attempt
via linearisation of $\xi^{-5/6}$ (Doc.~368, Sect.~5) fails with $23\,\%$
discrepancy, because $\xi^{-5/6}$ expands in $\xi^{1/6}$, not in $\xi$.

\paragraph{What distinguishes Route 2.}
The Koide route does not need this correction separately: it is contained in
$\cos(2/9)$. That is the reason for its precision -- and the reason it is not
algebraic. Routes 1, 3, 4 attempt to express the same result rationally or in
$\phi$-powers; the transcendental remainder is the correction they then have to
add back.

\section{Route 5: The Koide Scale from the T0 Ladder}

The Koide formula yields ratios; the scale $M$ remains external. Can it come
from Route~1?

\paragraph{What $M$ is.}
With $\sqrt{m_k} = \sqrt{M}\,a_k$ and $\sum_k a_k^2 = 6$ one has exactly
$M = (m_e+m_\mu+m_\tau)/6 = 313{.}850$ MeV \BM. This is normalisation, not
prediction. Equally exact: $M = m_\tau/a_\tau^2$ with $a_\tau^2 = 5{.}662$.

\paragraph{Finding.}
Inserting the bare masses from Route~1:
\[
  M_\text{bare} = \frac{v}{6}\Bigl(\tfrac{4}{3}\xi^{3/2} + \tfrac{16}{5}\xi
  + \tfrac{25}{9}\xi^{2/3}\Bigr) = 314{.}82\;\text{MeV}, \qquad
  \text{dev.\ } 0{.}31\,\% \quad \KM
\]
Route~5 is Route~1 summed. The correction is $+23{.}1\,\xi$ and is
arithmetically the weighted sum of the three individual corrections:
$-0{.}02\,\xi$ ($e$) $-2{.}78\,\xi$ ($\mu$) $+25{.}9\,\xi$ ($\tau$).
It is $\tau$-dominated because $M$ is $94\,\%$ $m_\tau/6$.

\paragraph{What this means.}
$M$ is not an independent scale. It follows from $\xi$ and $v$ with the
accuracy of Route~1 -- and with its open edge. A Galois-native single form
$M = c\,\xi^p\,v$ does not exist: $M$ is a sum of three different $\xi$-powers
($3/2$, $1$, $2/3$) that do not collapse into one (effective exponent
$0{.}747$, no rational coefficient for $p\in\{2/3,3/4,5/6,1\}$).

\paragraph{Warning against number games.}
Expressions such as $M \approx 81\sqrt{15}$ MeV ($0{.}044\,\%$) or
$M \approx \tfrac{7}{40}\,\alpha\,v$ ($0{.}1\,\%$) fit well but carry no
mechanism. With the building blocks $\{2,3,5,\phi,\sqrt2,\sqrt3,\sqrt5\}$
and exponents $-4\ldots4$, twelve combinations hit $M/E_0$ to within
$0{.}1\,\%$; $81\sqrt{15}$ is one of them, four others are better. The factor
$5/2$ in it appears in none of Docs.~292, 293, 338, 352, 367, 368. Such hits
do not belong in the corpus.

\section{What the Corpus Shows}

Read correctly, the five routes tell a single story -- and it is not about
leptons but about what FFGFT predicts at all.

FFGFT computes in natural units. The fundamental parameter $\xi = 4/30000$ is
dimensionless, the winding numbers are integers, the Galois orders are integers,
the icosahedral weights are algebraic numbers. None of this carries a unit. What
follows from this structure is therefore \emph{ratios}: $m_\mu/m_e$,
$m_\tau/m_e$, $Q = 2/3$, $1/\alpha$. And it is precisely these ratios that the
corpus hits, wherever it computes parameter-free, to $10^{-3}$ percent -- within
experimental uncertainty. The Koide route needs one matrix element of the $A_5$
geometry and nothing else. That is the theory's actual predictive power: the
\emph{shape} of the mass spectrum, not its size.

As soon as one demands a number in MeV, a scale must come from outside. FFGFT
takes a single declared anchor for this -- here $v = 246$ GeV, elsewhere $m_e$
or $H_0$ -- and this anchor serves conversion, not derivation. With the anchor
one inherits the percent class: the T0 ladder yields absolute masses to
$0{.}4$--$1{.}6$ percent, and the Koide scale $M$ is, as Route~5 shows, nothing
but this ladder summed. The two accuracy classes of the synopsis are thus not
two quality levels of the theory. They are the boundary between what FFGFT says
of its own accord (ratios, $10^{-3}\,\%$) and what it needs an external unit
for (absolute values, percent).

The open edge lies exactly on this boundary. The fractal-recursive correction --
$117\,\xi$ here, $27\xi/12$ there, hidden in $\cos(2/9)$ for Koide -- is the
amount by which the rational and algebraic approximations deviate from the
transcendental exact value. \emph{Measuring} it is easy; \emph{deriving} it from
the winding numbers is what is still missing. Until then: FFGFT predicts the
ratios of the lepton masses parameter-free and experimentally sharp, and their
absolute values with one anchor at percent level. The one is theory; the other
is theory plus a convention.
